\documentclass[11pt]{article}
\usepackage[margin=1in]{geometry}
\usepackage{amsmath,amssymb,amsthm,hyperref}
\newtheorem{theorem}{Theorem}[section]
\newtheorem{lemma}[theorem]{Lemma}
\newtheorem{corollary}[theorem]{Corollary}
\newcommand{\E}{\mathbb E}
\newcommand{\F}{\mathcal F}
\newcommand{\R}{\mathbb R}
\title{Conditional descent and almost-sure convergence of stochastic root iterations}
\author{}\date{}
\begin{document}\maketitle
\begin{abstract}
A conditional descent inequality with pathwise summable multiplicative and
additive errors yields a finite almost-sure limit and finite total restoring
work. We prove the result using predictable stopping, without requiring an
integrable infinite error budget. Applied to an initialized stochastic root
iteration, it gives almost-sure convergence and boundedness under global
restoration and square-summable, nonsummable steps. This explains a direct
connection between the root-energy argument and the later
almost-supermartingale approach to stochastic approximation.
\end{abstract}

\section{The conditional descent principle}
Throughout, $(\Omega,\F,\mathbb P)$ is a probability space with a filtration
$(\F_n)_{n\ge0}$. Conditional inequalities are understood almost surely.
All processes in the first two sections are real valued. A process is adapted
when its value at age $n$ is $\F_n$-measurable. Empty sums are zero and empty
products are one.

\begin{theorem}[Conditional descent]\label{thm:descent}
Let $X,A,B,C$ be adapted processes. Suppose that $X_n,B_n,C_n$ are integrable
for every finite $n$, that $X_n,A_n,B_n\ge0$ almost surely, and that
\begin{equation}\label{eq:row}
 \E[X_{n+1}\mid\F_n]\le(1+A_n)X_n-B_n+C_n.
\end{equation}
Assume also
\[
 \sum_{n=0}^\infty A_n<\infty,
 \qquad \sum_{n=0}^\infty |C_n|<\infty
 \quad\text{almost surely}.
\]
Then $X_n$ converges almost surely to a finite nonnegative random variable,
and $\sum_n B_n<\infty$ almost surely.
\end{theorem}
The additive errors may be signed. The assumptions impose no integrability
condition on $A_n$, on $A_nX_n$, or on either infinite error sum.
Consequently, a direct expectation bound for the full unnormalized right
side of \eqref{eq:row} need not be available.

The proof has two steps. Division by the accumulated multiplicative factors
produces an additive descent inequality whose terms remain integrable.
Stopping just before the predictable error prefix crosses a deterministic
level then permits supermartingale convergence. The stopping argument and
normalization are proved in Section~\ref{sec:proof}.

\begin{corollary}[Finite normalized account]\label{cor:finite}
Under the finite-age assumptions of Theorem~\ref{thm:descent}, set
\[
 P_0=1,\qquad P_n=\prod_{j=0}^{n-1}(1+A_j).
\]
For every integer $N\ge0$,
\begin{equation}\label{eq:finite}
 \E\frac{X_N}{P_N}
 +\sum_{n=0}^{N-1}\E\frac{B_n}{P_{n+1}}
 \le \E X_0+\sum_{n=0}^{N-1}\E\frac{C_n}{P_{n+1}}.
\end{equation}
No infinite summability assumption is needed for this finite inequality.
\end{corollary}
The terminal quantity in \eqref{eq:finite} is the actual normalized value at
age $N$. Each displayed expectation is finite because $P_n\ge1$. At $N=0$
the inequality is equality.

\section{Proof by predictable stopping}\label{sec:proof}
We use the nonnegative supermartingale convergence theorem: an integrable
nonnegative supermartingale converges almost surely to a finite limit. Its
$L^1$ norms are bounded by its initial expectation. We also use discrete
optional stopping for the process stopped at $n\wedge\tau$, which is a
bounded stopping time for each fixed $n$.

\begin{lemma}[A predictable lower envelope]\label{lem:localization}
Let $M$ be an integrable supermartingale. Suppose $Q_0=0$, that $Q_{n+1}$ is
$\F_n$-measurable, and that almost surely
\[
 M_n+Q_n\ge0\quad(n\ge0),\qquad \sup_{n\ge0}Q_n<\infty.
\]
Then $M_n$ converges almost surely to a finite limit.
\end{lemma}
\begin{proof}
For each positive integer $K$, define
\[
 \tau_K=\inf\{n\ge0:Q_{n+1}>K\},\qquad \inf\varnothing=\infty.
\]
This is a stopping time because $Q_{n+1}$ is known at age $n$.
For every $n$, the retained prefix satisfies
\begin{equation}\label{eq:noovershoot}
 Q_{n\wedge\tau_K}\le K.
\end{equation}
Indeed, this is immediate if $n\wedge\tau_K=0$. Otherwise, if the retained
index is $j\ge1$, then $j-1<\tau_K$, so $Q_j\le K$ by the definition of the
first crossing. This argument requires no monotonicity of $Q$.

The stopped process $M_{n\wedge\tau_K}$ is an integrable supermartingale,
and \eqref{eq:noovershoot} gives
\[
 M_{n\wedge\tau_K}+K\ge M_{n\wedge\tau_K}+Q_{n\wedge\tau_K}\ge0.
\]
It therefore converges to a finite limit almost surely. Intersect these
probability-one convergence events over all positive integers $K$ and with
the event on which $Q$ is bounded above. On any resulting path, choose
$K>\sup_n Q_n$. Then $\tau_K=\infty$, and the stopped sequence is the
original sequence for every age. This proves the assertion.
\end{proof}

\begin{lemma}[Additive descent]\label{lem:additive}
Let $Z,d,c$ be adapted and integrable at each finite age. Suppose
$Z_n,d_n\ge0$ almost surely, $\sum_n|c_n|<\infty$ almost surely, and
\[
 \E[Z_{n+1}\mid\F_n]\le Z_n-d_n+c_n.
\]
Then $\sum_n d_n<\infty$ almost surely and $Z_n$ has a finite nonnegative
almost-sure limit.
\end{lemma}
\begin{proof}
Write
\[
 D_n=\sum_{j<n}d_j,\qquad Q_n=\sum_{j<n}c_j,\qquad M_n=Z_n+D_n-Q_n.
\]
The conditional inequality makes $M$ a supermartingale. All its finite-age
terms are integrable. Moreover, $M_n+Q_n=Z_n+D_n\ge0$,
$Q_{n+1}$ is $\F_n$-measurable, and $Q_n$ converges almost surely by absolute
summability. Lemma~\ref{lem:localization} shows that $M_n$ has a finite limit.
On this probability-one event,
\[
 0\le D_n\le M_n+Q_n.
\]
The right side converges, so the increasing sequence $D_n$ is bounded and
has a finite limit. Hence $\sum_n d_n<\infty$, and
$Z_n=M_n-D_n+Q_n$ converges. Its limit is nonnegative.
\end{proof}

\begin{proof}[Proof of Theorem~\ref{thm:descent} and Corollary~\ref{cor:finite}]
To define denominators everywhere, use
$P_n=\prod_{j<n}(1+\max\{A_j,0\})$. On the simultaneous probability-one
event where all $A_j$ are nonnegative this agrees with the displayed product
in Corollary~\ref{cor:finite}. In particular, $P_n\ge1$ everywhere and
$P_{n+1}$ is $\F_n$-measurable. Set
\[
 Z_n=\frac{X_n}{P_n},\qquad
 d_n=\frac{B_n}{P_{n+1}},\qquad
 c_n=\frac{C_n}{P_{n+1}}.
\]
These variables are adapted and integrable, since their absolute values are
bounded by those of $X_n,B_n,C_n$, respectively. The predictable factor
$1/P_{n+1}$ is bounded by one, so it may be pulled through conditional
expectation. Dividing \eqref{eq:row} therefore gives
\begin{equation}\label{eq:normal}
 \E[Z_{n+1}\mid\F_n]\le Z_n-d_n+c_n.
\end{equation}
This calculation integrates only the normalized quantities; it never requires
integrability of $A_nX_n$. Taking expectations in \eqref{eq:normal} and
summing proves \eqref{eq:finite}.

For the infinite conclusion, $\sum_n|c_n|\le\sum_n|C_n|<\infty$ almost surely.
Lemma~\ref{lem:additive} gives convergence of $Z_n$ and summability of $d_n$.
On a path where $\sum_n A_n<\infty$, the product is increasing and satisfies
\[
 1\le P_n\le\exp\!\left(\sum_{j=0}^\infty A_j\right)<\infty.
\]
Thus $P_n\to P_\infty<\infty$, and $X_n=P_nZ_n$ converges to a finite
nonnegative limit. Finally,
\[
 \sum_{n=0}^\infty B_n
 =\sum_{n=0}^\infty P_{n+1}d_n
 \le P_\infty\sum_{n=0}^\infty d_n<\infty.
\]
The bound by $P_\infty$ is pathwise; no expectation of that bound is taken.
\end{proof}

\section{The initialized root iteration}
Let $E$ be a real Hilbert space, $x_*\in E$, and $F:E\to E$ continuous.
Consider the actual recursion
\begin{equation}\label{eq:root}
 X_{n+1}=X_n-a_nF(X_n)-a_n\xi_{n+1},\qquad n\ge0,
\end{equation}
with deterministic step sizes. All vector random variables below are strongly
measurable with respect to the indicated sigma algebra. Conditional
expectations of nonnegative variables may initially be interpreted in the
extended sense.

\begin{theorem}[Almost-sure root convergence]\label{thm:root}
Suppose there are constants $c>0$, $L\ge0$, $s\ge0$ such that, for every $x\in E$,
\begin{equation}\label{eq:geometry}
 \langle x-x_*,F(x)\rangle\ge c\|x-x_*\|^2,
 \qquad \|F(x)\|\le L\|x-x_*\|.
\end{equation}
Assume $X_0$ is $\F_0$-measurable and square integrable; $\xi_{n+1}$ is
$\F_{n+1}$-measurable; and
\[
 \E[\xi_{n+1}\mid\F_n]=0,\qquad
 \E[\|\xi_{n+1}\|^2\mid\F_n]\le s.
\]
Finally, suppose
\[
 a_n\ge0,\qquad \sum_n a_n=\infty,\qquad \sum_n a_n^2<\infty.
\]
Then $X_n\to x_*$ almost surely. In particular, $\sup_n\|X_n\|<\infty$
almost surely.
\end{theorem}
\begin{proof}
The extended conditional second-moment bound and the tower property give
$\E\|\xi_{n+1}\|^2\le s$. Thus each innovation is square integrable before
its vector conditional expectation is used. Continuity of $F$ and
\eqref{eq:root} give adaptation inductively. The growth bound in
\eqref{eq:geometry}, the triangle inequality and square integrability of
the innovations give square integrability of $X_n$ at every finite age.

Put $Z_n=\|X_n-x_*\|^2$. In the conditional square expansion, the cross term
with $\xi_{n+1}$ vanishes because $X_n-x_*-a_nF(X_n)$ is predictable.
The cross term is integrable by Cauchy--Schwarz. Therefore
\begin{align*}
 \E[Z_{n+1}\mid\F_n]
 &=\|X_n-x_*-a_nF(X_n)\|^2
   +a_n^2\E[\|\xi_{n+1}\|^2\mid\F_n]\\
 &\le(1+L^2a_n^2)Z_n-2ca_nZ_n+sa_n^2.
\end{align*}
Apply Theorem~\ref{thm:descent} with
\[
 A_n=L^2a_n^2,\qquad B_n=2ca_nZ_n,\qquad C_n=sa_n^2.
\]
All finite-age integrability assumptions have just been proved. The two error
series are finite, so $Z_n$ has a finite nonnegative limit and
$\sum_n a_n Z_n<\infty$ almost surely. A positive limit would imply
$Z_n\ge\varepsilon>0$ eventually and hence
$\sum_n a_nZ_n\ge\varepsilon\sum_{n\ \mathrm{large}}a_n=\infty$.
The limit must therefore be zero, which is equivalent to $X_n\to x_*$.
Every convergent sequence is bounded.
\end{proof}
No independence of the innovations is required. Their conditional centering
is the property used in the square expansion. There is also no restriction
on finitely many early steps, beyond their nonnegativity.

\begin{corollary}[Finite coordinate representations]\label{cor:coordinates}
Under Theorem~\ref{thm:root}, let $I$ be finite and let
$T:E\to\R^I$ be continuous and linear. Almost surely,
\[
 T X_n\longrightarrow T x_*,\qquad
 \sup_{n\ge0}\|T X_n\|_\infty<\infty.
\]
\end{corollary}
\begin{proof}
Apply continuity of $T$ on each convergent path, then boundedness of a
convergent sequence in the finite-coordinate maximum norm.
\end{proof}

\section{Connection with differential-equation methods}
The root-energy calculation starts from the initialized iteration and a
squared distance. The almost-supermartingale principle converts its
conditional descent into a pathwise conclusion. The lineage runs from the
stochastic root problem of Robbins and Monro~\cite{rm} to the conditional
convergence principle of Robbins and Siegmund~\cite{rs}; the modern proof analysis of Neri and
Powell uses the same product normalization and predictable stopping
\cite{np}. The Hilbert-space result above is a modern application of that
principle, rather than a claim about the exact scope of the original scalar
root theorem.

Bounded iterates are a substantive input in differential-equation tracking
arguments. Corollary~\ref{cor:coordinates} supplies that input for this
globally restoring subclass, using the same initialized trajectory.
A differential-equation application must still identify the represented
recursion and check its drift regularity, noise and remainder conditions.
Global restoration does not establish boundedness for arbitrary stochastic
approximation schemes. Nor does the argument above provide a finite
convergence rate, a general limit-set theorem, or a central limit theorem.

\begin{thebibliography}{9}
\bibitem{rm} H. Robbins and S. Monro.
\emph{A stochastic approximation method}.
The Annals of Mathematical Statistics 22(3), 400--407, 1951.
\bibitem{rs} H. Robbins and D. Siegmund.
\emph{A convergence theorem for non negative almost supermartingales and some
applications}. In \emph{Optimizing Methods in Statistics}, Elsevier, 1971,
pp.\ 233--257.
\bibitem{np} M. Neri and T. Powell.
\emph{A quantitative Robbins--Siegmund theorem}.
arXiv:2410.15986, version 2. In particular, the proof of the quantitative
Robbins--Siegmund theorem gives the product normalization and stopping
construction used here.
\end{thebibliography}
\end{document}
